%==========================================================
%  CHAPTER 4: Housing Affordability and Real Estate Markets
%==========================================================

\chapter{Housing Affordability and Real Estate Markets}
\label{ch:housing}

\begin{learningobjectives}
\item Explain why housing functions simultaneously as a consumption good and an investment asset, and describe the implications of this dual nature for market behaviour.
\item Apply the supply and demand model to housing markets and explain why price responses to demand shocks differ across cities.
\item Identify the key supply-side constraints---zoning, development charges, approval delays, and geography---that limit new housing construction in Canadian cities.
\item Analyze rental market dynamics and evaluate rent control as a policy response to rent inflation, drawing on theoretical and empirical evidence.
\item Describe homeownership gaps in Canada across generations, income groups, and regions, and connect these gaps to the price-to-income ratio.
\item Evaluate recent federal and provincial housing policy responses and assess their likely effectiveness.
\end{learningobjectives}

%----------------------------------------------------------
\section{Introduction: Canada's Housing Crisis}
%----------------------------------------------------------

In 2024, a 28-year-old graduate in Vancouver earning \$70,000 per year---a salary above the city's median---is priced out of the housing market. A one-bedroom condominium costs approximately \$700,000; a detached home with any outdoor space exceeds \$1.5 million. Saving a 20\% down payment at the rate of \$10,000 per year would take seven decades. Her parents, by contrast, purchased a similar home in 1990 for \$175,000---approximately twice their combined household income at the time.

This generational divide is the housing crisis in human terms. Economically, it represents a profound failure of markets and policy to keep housing supply pace with demand across Canadian cities over several decades. The consequences extend beyond individual hardship: when workers cannot afford to live in the cities where the most productive jobs are located, economic efficiency suffers. When housing wealth accumulates primarily among older, already-wealthy homeowners, inequality widens (as documented in Chapter~\ref{ch:inequality}). And when younger Canadians spend most of their income on rent, they save less, invest less in themselves, and have fewer children---all with long-run macroeconomic consequences.

Canada now has some of the least affordable housing among OECD nations. By 2024, the ratio of the average home price to the average household income exceeded 10:1 nationally and reached 12:1 in Vancouver and Toronto. Oxford Economics regularly ranks Vancouver and Toronto among the top five least affordable housing markets in the world.

%----------------------------------------------------------
\section{Housing as an Economic Good}
%----------------------------------------------------------

\subsection{The Dual Nature of Housing}

Housing occupies an unusual position in economics because it serves two distinct economic functions simultaneously:

\textbf{Housing as a consumption good.} Like food or clothing, housing provides direct utility---shelter, warmth, privacy, space for family life. Every household must consume some housing. For most Canadians, it is their largest monthly expenditure.

\textbf{Housing as an investment asset.} Owner-occupied housing also functions as a store of wealth. Homeowners benefit (or suffer) from capital gains as home prices change. They earn an \term{imputed rent}---the rental income they would receive if they rented their home to someone else, which they effectively receive by living in it rent-free.

\begin{defbox}{Imputed Rent}
\textbf{Imputed rent} is the estimated rental value of owner-occupied housing---the income an owner would receive by renting their property to a tenant, or equivalently, the rent they save by owning rather than renting.

\textit{Economic significance:} Imputed rent is counted in GDP (national accounts include it as a component of household income). It also means that owning a home that has appreciated substantially provides a flow of services as well as a capital gain.

\textit{Canadian example:} A homeowner in Toronto whose house appreciated from \$500,000 to \$1.2 million has gained \$700,000 in net worth. They also receive the equivalent of market-rate rent (\$3,000+/month) as a flow of housing services---a form of income not subject to income tax under the principal residence exemption.
\end{defbox}

This dual nature creates distinctive market dynamics:

\begin{itemize}
  \item Housing is heterogeneous---no two units are identical in location, age, condition, and amenities.
  \item Housing is immobile---supply in one location cannot move to serve demand in another.
  \item Housing transactions are infrequent and high-cost (land transfer taxes, legal fees, real estate commissions typically total 3--6\% of purchase price).
  \item Housing supply adjusts slowly---new construction takes 2--5 years from zoning approval to occupancy.
  \item Housing is heavily regulated---zoning, building codes, rent control, condo legislation, strata law.
  \item Expectations matter---if buyers expect prices to continue rising, speculative demand can persist even when prices are already very high.
\end{itemize}

\begin{keytakeaway}
The dual nature of housing---simultaneously a necessity and an investment---creates a fundamental tension in housing policy. Policies that make housing cheaper as a consumption good (e.g., increased supply, rent controls) may be opposed by existing homeowners for whom housing is their primary investment asset. This political economy dynamic helps explain why supply-constraining policies persist even when their affordability costs are well documented.
\end{keytakeaway}

%----------------------------------------------------------
\section{The Economics of Housing Supply and Demand}
%----------------------------------------------------------

\subsection{Housing Demand}

The demand for housing in a given city depends on several factors:

\begin{itemize}
  \item \textbf{Income.} Higher household incomes shift demand for housing outward, as households seek more space, better locations, and higher-quality units.
  \item \textbf{Population growth and household formation.} More people---through natural increase, immigration, or in-migration from other regions---increase the number of households seeking housing.
  \item \textbf{Mortgage interest rates.} Lower rates reduce the monthly cost of a given purchase price, increasing the quantity of housing demanded. This connection between interest rates and housing demand is a major reason why the Bank of Canada's 2022--2023 rate hike cycle dampened housing prices (see Chapter~\ref{ch:inflation}).
  \item \textbf{Expectations of price appreciation.} If buyers expect continued price increases, speculative demand rises. Expectations can become self-fulfilling in the short run.
  \item \textbf{Immigration and demographic composition.} Canada's immigration targets---reaching 500,000 new permanent residents per year---generate substantial additional housing demand. New arrivals initially concentrate in Vancouver, Toronto, and Montreal.
\end{itemize}

\subsection{Housing Supply and Its Inelasticity}

The supply of new housing depends on land availability, construction costs, development regulations, approval timelines, and the profitability expectations of builders. The critical insight for understanding Canadian housing markets is that housing supply is extremely \term{price inelastic}---it responds slowly and by small amounts to increases in housing prices.

\begin{defbox}{Price Elasticity of Supply}
The \textbf{price elasticity of supply} measures the percentage change in quantity supplied in response to a 1\% increase in price: $E_s = \frac{\% \Delta Q_s}{\% \Delta P}$.

\textit{Elastic supply} ($E_s > 1$): Quantity supplied responds strongly to price. Demand increases raise quantity more than price.

\textit{Inelastic supply} ($E_s < 1$): Quantity supplied responds weakly to price. Demand increases raise price more than quantity.

\textit{Canadian context:} Housing supply in Vancouver and Toronto is estimated to have price elasticities of approximately 0.2--0.4---highly inelastic. In comparison, Dallas and Houston have elasticities closer to 2.0, consistent with their much lower price growth despite strong population inflows.
\end{defbox}

\begin{figure}[H]
\centering
\begin{tikzpicture}
% Panel A: Elastic supply
\begin{scope}
\begin{axis}[
  at={(0,0)},
  width=7cm, height=7.5cm,
  xlabel={Quantity of Housing},
  ylabel={Price Level},
  xmin=0, xmax=10, ymin=0, ymax=10,
  xtick=\empty, ytick=\empty,
  axis line style={-Stealth, thick},
  axis lines=left,
  title={\small Panel A: Elastic Supply (Houston)},
  title style={color=chaptergray},
  clip=false
]
% Elastic supply (flatter)
\addplot[thick, color=policyblue, domain=2:9] {0.6*x + 1}
  node[right, font=\scriptsize] {$S$ (elastic)};
% D0
\addplot[thick, color=chaptergray, dashed, domain=1:9] {-x + 9}
  node[right, font=\scriptsize] {$D_0$};
% D1 (shifted right)
\addplot[thick, color=unbcgreen, domain=2:9.5] {-x + 11}
  node[right, font=\scriptsize] {$D_1$};
% Equilibrium points
\addplot[only marks, mark=*, color=black, mark size=2.5pt]
  coordinates {(4.91,3.95)};
\addplot[only marks, mark=*, color=policyblue, mark size=2.5pt]
  coordinates {(5.87,4.52)};
% Arrows showing small price rise, large quantity rise
\draw[<->, thick, color=dataorange] (axis cs:5.8,3.95) -- (axis cs:5.8,4.52)
  node[right, midway, font=\scriptsize, color=dataorange] {Small $\Delta P$};
\draw[<->, thick, color=unbcgreen] (axis cs:4.91,2.0) -- (axis cs:5.87,2.0)
  node[below, midway, font=\scriptsize, color=unbcgreen] {Large $\Delta Q$};
\end{axis}
\end{scope}
% Panel B: Inelastic supply
\begin{scope}
\begin{axis}[
  at={(7.5cm,0)},
  width=7cm, height=7.5cm,
  xlabel={Quantity of Housing},
  ylabel={Price Level},
  xmin=0, xmax=10, ymin=0, ymax=10,
  xtick=\empty, ytick=\empty,
  axis line style={-Stealth, thick},
  axis lines=left,
  title={\small Panel B: Inelastic Supply (Vancouver)},
  title style={color=chaptergray},
  clip=false
]
% Inelastic supply (steeper)
\addplot[thick, color=policyblue, domain=3.5:7] {2.5*x - 7}
  node[right, font=\scriptsize] {$S$ (inelastic)};
% D0
\addplot[thick, color=chaptergray, dashed, domain=1:9] {-x + 9}
  node[right, font=\scriptsize] {$D_0$};
% D1 (shifted right)
\addplot[thick, color=unbcgreen, domain=1:9] {-x + 11}
  node[right, font=\scriptsize] {$D_1$};
% Equilibrium points
\addplot[only marks, mark=*, color=black, mark size=2.5pt]
  coordinates {(4.57,4.43)};
\addplot[only marks, mark=*, color=policyblue, mark size=2.5pt]
  coordinates {(4.86,6.14)};
% Arrows
\draw[<->, thick, color=dataorange] (axis cs:5.5,4.43) -- (axis cs:5.5,6.14)
  node[right, midway, font=\scriptsize, color=dataorange] {Large $\Delta P$};
\draw[<->, thick, color=unbcgreen] (axis cs:4.57,2.0) -- (axis cs:4.86,2.0)
  node[below, midway, font=\scriptsize, color=unbcgreen] {Small $\Delta Q$};
\end{axis}
\end{scope}
\end{tikzpicture}
\caption{Demand Shocks Under Elastic vs.\ Inelastic Housing Supply. When demand increases
  (D$_0$ to D$_1$) in a market with elastic supply (Panel A), most of the adjustment occurs
  in quantity, with modest price increases. In an inelastically supplied market (Panel B,
  typical of Vancouver), the same demand shock produces a large price increase and a
  small quantity increase. The difference in supply elasticity is the fundamental reason
  why comparable demand shocks produce dramatically different price outcomes across cities.}
\label{fig:housing_sd}
\end{figure}

%----------------------------------------------------------
\section{Housing Supply Constraints in Canada}
%----------------------------------------------------------

\subsection{Zoning and Land Use Regulation}

The most pervasive supply constraint in Canadian cities is \term{exclusionary zoning}: municipal land-use regulations that prohibit multi-unit residential buildings (apartments, townhouses, duplexes) in most residential neighbourhoods, permitting only detached single-family homes.

In Vancouver, approximately 70\% of residential land was zoned exclusively for single-family houses before provincial intervention in 2023. In Toronto, over 60\% of land within the old city was single-family zoned. These regulations reflect the political preferences of existing homeowners, who benefit from restricted supply (which supports property values) and have strong incentives to participate in local planning processes.

The economics of exclusionary zoning: zoning restrictions act like a quantity control on the housing market, limiting the total supply below what would emerge in an unregulated market. The result is an implicit tax on new housing---the gap between the market price of a unit and its production cost---which can reach several hundred thousand dollars per unit in supply-constrained cities.

\begin{canexample}{BC's Housing Statutes Amendment Act (Bill 44, 2023)}
In November 2023, the British Columbia government passed Bill 44 --- the
\textit{Housing Statutes (Residential Development) Amendment Act} --- marking one of the
most significant zoning reforms in Canadian history. Key provisions:
\begin{itemize}
  \item All municipalities above a population threshold must permit \textbf{small-scale
        multi-family housing} (up to 4 units) on any lot that previously allowed a
        single-family home.
  \item Transit-oriented areas near SkyTrain and bus rapid transit stations must allow
        higher densities (up to 20 storeys within 200m of rapid transit).
  \item Municipalities cannot use design, parking, or character bylaws to effectively
        block projects that comply with the new rules.
\end{itemize}
The legislation was modelled partly on Auckland, New Zealand's 2016 upzoning, which research
found increased housing supply and meaningfully reduced rent growth relative to comparable
cities that did not reform. Early evidence from BC suggests development applications in
affected municipalities have increased.
\end{canexample}

\subsection{Development Charges and Approval Delays}

Development charges---fees levied by municipalities on developers to fund infrastructure (roads, water, parks) associated with new development---are a significant cost add-on to new construction. In Toronto, development charges per unit averaged approximately \$150,000 as of 2023---one of the highest levels in the world. These charges add directly to the cost of new housing and deter marginal projects.

Municipal approval timelines in major Canadian cities typically range from 2 to 4 years from application to approval. Each year of delay imposes carrying costs on developers (typically 6--8\% of land value per year at current interest rates), which are passed through to the final price of units. Streamlining approval processes is low-cost from a government fiscal perspective but politically difficult when approval processes involve extensive public consultation.

\subsection{Construction Labour Shortages}

The construction labour force in Canada is aging: approximately 22\% of trades workers were over 55 in 2022, with large cohorts approaching retirement. Training new workers in skilled trades requires 4--5 years of apprenticeship, creating a medium-term supply constraint. Immigration of workers with construction credentials is one policy response, but credential recognition barriers have historically been significant.

\subsection{The CMHC Housing Shortfall}

\begin{databox}{Canada's Housing Supply Gap}
CMHC's 2022 Housing Supply Report estimated that Canada needs approximately \textbf{3.5 million additional housing units} by 2030 beyond what was already planned, to restore affordability to 2003--2004 levels.

\begin{itemize}
  \item Current annual construction pace: approximately 200,000--250,000 starts
  \item Required pace to meet the shortfall target: approximately 700,000+ per year
  \item BC shortfall: approximately 570,000 units
  \item Ontario shortfall: approximately 1.5 million units
  \item Prince George and surrounding area: estimated shortfall of several thousand units
\end{itemize}

At current construction rates, the target will not be met.
\source{CMHC, Housing Supply Report, 2022; CMHC Housing Market Outlook.}
\end{databox}

%----------------------------------------------------------
\section{Rent Inflation and the Rental Market}
%----------------------------------------------------------

\subsection{Rental Market Dynamics}

Canada's rental market tightened dramatically between 2021 and 2024. National rental vacancy rates fell to approximately 1.5\%---well below the approximately 3\% rate considered ``balanced'' between landlord and tenant power. Rents surged across the country:

\begin{figure}[H]
\centering
\begin{tikzpicture}
\begin{axis}[
  xbar,
  bar width=0.5cm,
  width=12cm, height=8cm,
  symbolic y coords={Winnipeg, Prince George, Halifax, Montreal, Ottawa, Calgary, Toronto, Vancouver},
  ytick=data,
  y tick label style={font=\small},
  xlabel={Average 2-Bedroom Rent (CAD/month, 2023)},
  xlabel style={font=\small},
  xmin=0, xmax=3800,
  xtick={0,500,1000,1500,2000,2500,3000,3500},
  xticklabel style={font=\small},
  xmajorgrids=true, grid style={dashed, gray!25},
  every axis plot/.append style={fill=exampleteal!60, draw=exampleteal}
]
\addplot coordinates {
  (1450,Winnipeg)(1500,Prince George)(1800,Halifax)(1800,Montreal)
  (2100,Ottawa)(2200,Calgary)(2800,Toronto)(3200,Vancouver)
};
% National average line
\addplot[thick, dashed, color=dataorange] coordinates {(2050,Winnipeg)(2050,Vancouver)}
  node[above right, font=\scriptsize, color=dataorange] {National avg. $\approx\$2{,}050$};
\end{axis}
\end{tikzpicture}
\caption{Average Two-Bedroom Rent in Selected Canadian Cities, 2023. Vancouver and Toronto
  rents are roughly double those in more affordable markets. Prince George remains
  substantially more affordable than major metropolitan centres, though local rents
  have risen sharply since 2021.}
\label{fig:rents}
\source{CMHC Rental Market Report; Rentals.ca National Rent Report, 2023. Estimates.}
\end{figure}

\subsection{Rent Control: Theory and Evidence}

\begin{defbox}{Rent Control}
\textbf{Rent control} refers to government regulations that limit the rent a landlord may charge for a residential unit. It may take the form of an absolute cap on rent, or a cap on the permissible annual rent increase.

\textit{Simple explanation:} The government prevents landlords from charging more than a specified amount.

\textit{Economic intuition:} If the controlled rent is set below the market equilibrium rent, a shortage emerges---quantity demanded exceeds quantity supplied at the controlled price.
\end{defbox}

\textbf{Standard economic analysis of rent control} predicts:
\begin{enumerate}
  \item \textbf{Shortage:} At a controlled rent below equilibrium, more tenants want units than are available. Vacancy falls; waiting lists emerge; queuing and informal rationing replace price rationing.
  \item \textbf{Quality deterioration:} Landlords facing rent caps have reduced incentive to maintain or improve properties.
  \item \textbf{Supply reduction:} Below-market returns on rental investment reduce new rental construction and cause existing landlords to convert units to condominiums, strata-titled buildings, or other uses.
  \item \textbf{Misallocation:} Tenants in rent-controlled units have reduced incentive to move even when their circumstances change (household size, employment location), leading to inefficient matching of households to units.
\end{enumerate}

\textbf{Empirical evidence} broadly supports the long-run supply effects. Diamond, McQuade, and Qian (2019) studied a San Francisco ballot initiative that expanded rent control and found it reduced the supply of rental housing by approximately 15\%, as landlords converted properties or redeveloped them to avoid controls. While existing rent-controlled tenants benefited, the overall rental market tightened, driving up rents for non-controlled units.

\textbf{Canadian context:} Rent regulation varies by province:
\begin{itemize}
  \item \textbf{Ontario:} Rent control applies to most units built before November 2018. Annual increases capped at the Ontario Rent Increase Guideline (tied to inflation). Units built after 2018 are exempt to encourage new supply.
  \item \textbf{British Columbia:} Annual rent increase cap tied to inflation (currently BC CPI+2\%). No exemption for new units (unlike Ontario).
  \item \textbf{Quebec:} The Tribunal administratif du logement (TAL) sets guidelines for rent increases and adjudicates disputes.
  \item \textbf{Alberta, Saskatchewan, New Brunswick:} No rent control.
\end{itemize}

\begin{thinkbox}
A city council is considering extending Ontario-style rent control to all new apartment buildings (currently exempt for units built after 2018), arguing it is unfair that existing tenants are protected but new tenants in new buildings are not. Using economic theory and the empirical evidence discussed above, analyze the likely short-run and long-run effects of this policy. Who would benefit? Who would be harmed? What does the evidence from Diamond et al.\ suggest about the overall effect on renters as a group?
\end{thinkbox}

%----------------------------------------------------------
\section{Homeownership Gaps in Canada}
%----------------------------------------------------------

\subsection{The Generational Divide}

Canada's homeownership rate has diverged dramatically across birth cohorts. Overall, approximately 66\% of Canadian households own their home. But this aggregate masks a profound generational split:

\begin{figure}[H]
\centering
\begin{tikzpicture}
\begin{axis}[
  ybar,
  bar width=0.4cm,
  width=13cm, height=7.5cm,
  symbolic x coords={Under 35, 35--44, 45--54, 55--64, 65+},
  xtick=data,
  x tick label style={font=\small},
  ylabel={Homeownership Rate (\%)},
  ylabel style={font=\small},
  ymin=0, ymax=90,
  ytick={0,20,40,60,80},
  yticklabel style={font=\small},
  ymajorgrids=true, grid style={dashed, gray!25},
  enlarge x limits=0.12,
  legend style={at={(0.98,0.98)}, anchor=north east, font=\small},
  legend entries={1981, 2021}
]
\addplot[fill=unbcgreen!40, draw=unbcgreen] coordinates {
  (Under 35, 42)(35--44, 62)(45--54, 70)(55--64, 72)(65+, 72)
};
\addplot[fill=policyblue!60, draw=policyblue] coordinates {
  (Under 35, 36)(35--44, 55)(45--54, 67)(55--64, 74)(65+, 78)
};
\end{axis}
\end{tikzpicture}
\caption{Homeownership Rate by Age Group, Canada, 1981 vs.\ 2021. Homeownership rates
  among Canadians under 35 and aged 35--44 have declined since 1981, while rates among
  seniors have increased. Young Canadians today are less likely to own homes than their
  parents were at the same age, driven primarily by the deterioration of price-to-income
  ratios in major urban centres.}
\label{fig:homeownership_age}
\source{Statistics Canada, Census of Population.}
\end{figure}

The decline in young homeownership is not driven by changing preferences. Surveys consistently show that Millennials and Gen Z want to own at rates comparable to previous generations. The barrier is the price-to-income ratio. In 1976, the average Toronto home cost approximately 2$\times$ the average household income. By 2022, the ratio exceeded 12$\times$. A first-time buyer purchasing the national benchmark home (\$730,000) with a 10\% down payment (\$73,000) requires a household income of approximately \$140,000 to qualify for the mortgage under current stress-test rules---a threshold out of reach for most Canadian families.

\subsection{Income and Regional Gaps}

Homeownership is strongly correlated with income. The top quintile of households has a homeownership rate of approximately 85\%; the bottom quintile, approximately 25\%. Housing serves as both a consumption good and a vehicle for wealth accumulation---the two functions compound each other for those who can enter the market, and exclude those who cannot.

Regional disparities are large: homeownership remains relatively affordable in Prairie cities (Winnipeg benchmark price approximately \$370,000) and in secondary markets like Prince George (approximately \$450,000) compared to Vancouver (approximately \$1.2 million) and Toronto (approximately \$1.05 million). This relative affordability explains much of the internal migration toward Prairie and northern cities observed since 2020.

\subsection{Racial and Immigrant Homeownership Gaps}

Racialized Canadians and recent immigrants face lower homeownership rates than non-racialized Canadians, even after controlling for income:

\begin{itemize}
  \item White Canadians: approximately 73\% homeownership
  \item South Asian Canadians: approximately 63\%
  \item Chinese Canadians: approximately 68\%
  \item Black Canadians: approximately 47\%
  \item Recent immigrants (arrived within 5 years): approximately 28\%
\end{itemize}

These gaps reflect multiple factors: lower median incomes (reducing purchasing power), discrimination in mortgage lending, lower family wealth (limiting down payment support), and concentration in the most expensive rental markets (where savings are harder to accumulate).

%----------------------------------------------------------
\section{Housing Wealth Inequality}
%----------------------------------------------------------

Home equity is the largest single component of wealth for most Canadian households. The dramatic appreciation of Canadian home prices since 2000 has transferred enormous wealth to existing homeowners---disproportionately older, higher-income, non-racialized Canadians---while imposing rising costs on renters and first-time buyers.

\begin{canexample}{The Housing Boom as a Wealth Transfer}
A homeowner who purchased a detached home in Toronto for \$300,000 in 2005 and sold in
2022 for approximately \$1.1 million realized a \textbf{capital gain of \$800,000}---fully
tax-free under Canada's principal residence exemption. This gain equates to approximately
15 years of the average Canadian pre-tax salary.

Over the same period, a renter in an equivalent property paid approximately \$1,500/month
in 2005 rising to \$2,800/month by 2022---approximately \$430,000 in cumulative rent
payments, with no asset to show at the end. The renter accumulated no housing wealth;
the homeowner accumulated \$800,000.

This is not a story of individual virtue or failure. It is a consequence of a market
structure in which housing functions as both shelter and investment, and in which an
undersupply of housing inflated values beyond any productivity-related justification.
The distributional consequence: housing has become the primary mechanism through which
wealth inequality between and within generations has widened in Canada.
\end{canexample}

\textbf{The ``Bank of Mum and Dad.''} A CIBC report found that approximately 30\% of first-time homebuyers in Canada received gifts or loans from family members, averaging approximately \$130,000 per recipient. This creates a two-tier market: Canadians with wealthy parents can access homeownership (and the wealth accumulation it provides); those without cannot. Parental wealth transfers thus reinforce existing inequality across generations---a phenomenon connecting directly to the discussion of inherited wealth in Chapter~\ref{ch:inequality}.

%----------------------------------------------------------
\section{Northern BC and Prince George: A Regional Perspective}
%----------------------------------------------------------

While national headlines focus on Vancouver and Toronto, Northern BC's housing dynamics are distinctive and worth examining as a regional case study.

\textbf{Resource economy volatility.} Prince George's economy is tied to the forestry and resource sectors. Sawmill closures, lumber price downturns, or pulp mill shutdowns can rapidly reduce housing demand, while mine openings or LNG-related construction can trigger rapid demand surges. This boom-bust pattern creates particular challenges for developers and long-term homeowners whose equity is tied to commodity cycles.

\textbf{Relative affordability.} Prince George's benchmark home price of approximately \$450,000 (2024) is dramatically below Vancouver (\$1.2 million) and Toronto (\$1.05 million). This relative affordability has attracted internal migrants---particularly remote workers and retirees from the Lower Mainland---since 2020. While this increases local wealth and spending, it also competes with local buyers and can drive prices up for residents whose incomes are calibrated to the local labour market.

\textbf{Tightening rental market.} Despite lower ownership prices, Prince George's rental market has tightened significantly since 2021. Vacancy rates fell below 2\% in 2022--2023, and average rents for a two-bedroom unit rose from approximately \$1,100 in 2019 to approximately \$1,500--\$1,600 by 2024. For service-sector workers, students, and Indigenous urban residents---groups with lower average incomes---this tightening has created genuine affordability stress.

\textbf{Indigenous housing in Northern BC.} Urban Indigenous residents in Prince George face particularly acute housing challenges. They are overrepresented in unaffordable, overcrowded, or inadequate housing, and face barriers to homeownership including lower median incomes, limited access to traditional mortgage financing, and discrimination. On-reserve housing in surrounding First Nations communities faces chronic shortfalls in quantity and quality---a product of decades of federal under-investment and the jurisdictional barriers created by the Indian Act.

%----------------------------------------------------------
\section{Federal, Provincial, and Municipal Housing Policy}
%----------------------------------------------------------

\subsection{Recent Federal Initiatives}

\textbf{First Home Savings Account (FHSA), 2023.} A new registered savings account allowing first-time buyers to save up to \$8,000 per year (lifetime limit \$40,000), with tax-deductible contributions and tax-free withdrawals for home purchase. It combines features of the RRSP (contribution deductibility) and TFSA (tax-free withdrawals). Effective for those with savings capacity; of limited help to households who cannot save.

\textbf{Housing Accelerator Fund.} A \$4 billion federal program providing grants to municipalities that commit to meaningful zoning reform---particularly in increasing density near transit. Modelled on housing supply programs in the United States, with conditionality on reform commitments. Cities including Hamilton, London, and Halifax have signed agreements and committed to allowing more missing middle housing.

\textbf{Foreign Buyer Ban.} A 2-year ban (2023--2025) on most non-resident foreign nationals purchasing Canadian residential property. The policy rationale was that foreign investment was contributing to price inflation. The empirical evidence is mixed: foreign buyers represent a relatively small share of transactions in most markets (CMHC estimated approximately 3--5\% of total transactions even in peak years). The ban has had limited effect on affordability, reinforcing the view that foreign demand is not the primary driver of Canada's housing crisis.

\textbf{Apartment Construction Loan Program.} Low-cost federal loans for purpose-built rental apartment construction, aimed at incentivizing the rental supply that the private market has underproduced.

\subsection{Provincial and Municipal Approaches}

\textbf{British Columbia:} In addition to Bill 44 (discussed above), BC has introduced transit-oriented zoning (mandating higher densities near rapid transit), short-term rental regulations (restricting Airbnb to primary residences in most municipalities), and a speculation and vacancy tax on properties left unoccupied.

\textbf{Ontario:} The \textit{More Homes Built Faster Act} (2022) attempted to accelerate approvals, reduce development charges, and allow more as-of-right development. Implementation has been contested, with some municipalities challenging provincial overrides of local zoning.

\textbf{Quebec:} A stronger tradition of non-profit and social housing has meant lower levels of housing unaffordability relative to Toronto and Vancouver, though Montreal has seen significant rent increases since 2020.

%----------------------------------------------------------
\section{Policy Debate}
%----------------------------------------------------------

\begin{policydebate}{Is Zoning Reform the Solution to Canada's Housing Crisis?}

\textbf{The Core Claim}

A growing consensus among economists holds that restrictive zoning---particularly single-family exclusionary zoning---is the primary cause of Canada's housing affordability crisis. The logic is straightforward: if demand consistently exceeds supply because supply is artificially constrained by regulation, prices rise. Remove the constraints, supply increases, and price growth moderates.

\vspace{8pt}
\textbf{Side A: Yes---Zoning Reform is the Key Solution}

\begin{itemize}
  \item Academic consensus: Research by Glaeser and Gyourko, Hsieh and Moretti, Albouy and Ehrlich, and others documents that restrictive zoning is the primary driver of housing cost divergence across cities globally.
  \item Tokyo as a counter-example: Japan's land-use system allows near-uniform density across the country. Tokyo---a city larger than the Greater Toronto Area---has maintained relatively affordable housing by permitting high-density development wherever demand exists.
  \item Auckland natural experiment: Auckland's 2016 upzoning permitting density across most residential land was followed by increased housing supply and slower rent growth than comparable New Zealand cities. A 2022 study found the reform created approximately 4\% more housing than would have otherwise been built.
  \item Houston comparison: Houston has no zoning by-law. It consistently maintains home prices at approximately 3$\times$ median income despite large population inflows, while similarly growing cities with restrictive zoning see price-to-income ratios exceed 10$\times$.
  \item Supply reform is fiscally cheap: Changing zoning by-laws costs essentially nothing from a government budget perspective.
\end{itemize}

\textbf{Side B: Supply Alone Is Insufficient}

\begin{itemize}
  \item New market-rate supply ``filters'' to lower-income households slowly, over 20--30 years. In the short to medium run (the relevant policy horizon), upzoning primarily produces luxury and market-rate units that do not serve low-income renters.
  \item Investor absorption: If speculative demand is sufficiently strong, new supply may be absorbed by investors rather than owner-occupiers, with limited near-term price effect.
  \item Community impacts: Rapid densification without infrastructure investment strains schools, transit, and parks. Displacement of lower-income residents from upzoned neighbourhoods is a documented concern.
  \item The social housing deficit: Market supply will never adequately house the bottom 10--15\% of the income distribution, for whom shelter subsidy and public housing are necessary.
\end{itemize}

\textbf{Evidence and Assessment}

The peer-reviewed economic literature strongly supports the supply-side argument for the long run. The critics are correct that market-rate supply does not immediately and fully address affordability for low-income renters---a different policy intervention is needed for that population. But restricting supply while blaming demand-side factors (speculation, foreign buyers) for high prices is not supported by the evidence.

\textbf{Synthesis}: Zoning reform is necessary but not sufficient. A comprehensive housing strategy requires: (1) zoning reform to allow density; (2) streamlined approvals and reduced development charges; (3) significant public and non-profit housing investment for low-income households; (4) tenant protections during the supply transition; and (5) infrastructure investment to support densification.
\end{policydebate}

%----------------------------------------------------------
\section*{Chapter Summary}
%----------------------------------------------------------

\begin{itemize}
  \item Housing has a \textbf{dual nature} as both a consumption good and an investment asset, creating distinctive market dynamics including speculative demand, slow supply adjustment, and political economy complications.
  \item Housing supply is \textbf{highly price inelastic} in most Canadian cities, meaning demand increases translate primarily into higher prices rather than greater quantities. This is the root cause of affordability deterioration.
  \item The primary supply-side constraint is \textbf{exclusionary single-family zoning}, which limits density in most Canadian residential neighbourhoods. Development charges, approval delays, and construction labour shortages compound this barrier.
  \item CMHC estimates Canada needs approximately \textbf{3.5 million additional housing units} by 2030 beyond current plans to restore affordability. Current construction rates are far below this pace.
  \item The \textbf{rental market} has tightened dramatically since 2021, with national vacancy rates below 2\%. Rent control can protect existing tenants but creates long-run supply disincentives, as documented by Diamond et al.\ for San Francisco.
  \item Canada has a pronounced \textbf{generational homeownership gap}: rates among Canadians under 35 have declined since 1981. The cause is not changing preferences but a deterioration of the price-to-income ratio from approximately 2:1 in the 1970s to 12:1 in major cities today.
  \item Housing appreciation has created a \textbf{housing wealth gap} between existing owners (who have accumulated large gains) and renters and non-owners (who have not), amplifying overall inequality.
  \item \textbf{Northern BC and Prince George} present a distinctive regional case: resource sector volatility, relative affordability compared to major metros, tightening rental markets since 2021, and acute Indigenous housing needs.
  \item Recent federal and provincial policy responses---the FHSA, Housing Accelerator Fund, BC Bill 44---represent meaningful steps but are unlikely to fully close the supply gap in the near term.
  \item The \textbf{evidence from comparative housing economics} (Tokyo, Houston, Auckland) strongly supports supply reform (particularly zoning liberalization) as necessary to restore affordability, though complementary policies for low-income households are also essential.
\end{itemize}

%----------------------------------------------------------
\section*{Review Questions}
%----------------------------------------------------------

\begin{enumerate}

\item \textbf{(Concept)} Explain why housing supply is more price-inelastic than the supply of most consumer goods. What are the main factors that prevent a rapid supply response to rising housing prices in Canadian cities?

\item \textbf{(Application)} Draw two supply-demand diagrams for housing: one representing a city with elastic supply (e.g., Calgary) and one with inelastic supply (e.g., Vancouver). Show the effect of a large increase in immigration-driven housing demand on each. For each city, indicate what happens to price and quantity. What explains the difference?

\item \textbf{(Concept)} Define rent control and explain the standard economic prediction for its effects. Under what conditions might rent control impose lower costs than the theory predicts?

\item \textbf{(Application)} A renter and a homeowner in the same Toronto neighbourhood have identical incomes of \$90,000 per year. Describe how their economic situations have diverged since 2005 in terms of (a) housing costs, (b) wealth accumulation, and (c) financial flexibility. What does this illustrate about housing as an investment asset?

\item \textbf{(Critical Thinking)} The federal government's Foreign Buyer Ban (2023--2025) targeted non-resident foreign nationals purchasing Canadian residential property. What economic hypothesis underlies this policy? What does the evidence suggest about how large a role foreign buyers play in Canada's housing prices? Does the evidence justify the policy?

\item \textbf{(Policy)} Evaluate the First Home Savings Account (FHSA) as a housing affordability policy. Who benefits directly from the FHSA? Does it address the demand side or the supply side of the housing market? What groups of first-time buyers are least helped by the FHSA?

\item \textbf{(Application)} How does housing in Prince George differ from housing in Vancouver, and what economic factors explain the differences? What are the implications of this difference for internal migration patterns?

\end{enumerate}

%----------------------------------------------------------
\section*{Discussion Questions}
%----------------------------------------------------------

\begin{enumerate}

\item Should adequate housing be guaranteed as a legal right in Canada? What would this mean in practice, and what would be the economic implications of such a guarantee?

\item Existing homeowners in Canada benefit financially when zoning is restrictive and housing supply is constrained, because their home values are higher than they would be in a freely supplied market. This creates a political conflict between existing homeowners (who want to preserve home values) and prospective buyers (who want lower prices). How should democratic governments navigate this conflict? Does the majority necessarily determine the right outcome?

\item The principal residence exemption allows Canadians to sell their primary home tax-free, even on gains of millions of dollars. Should this exemption be reformed or capped? What would be the economic effects of doing so?

\end{enumerate}

%----------------------------------------------------------
\section*{Exercises}
%----------------------------------------------------------

\begin{enumerate}

\item \textbf{Data Exercise.} Visit the CMHC website (\url{www.cmhc-schl.gc.ca}) and find the most recent Housing Market Outlook for Prince George (or the nearest published report for Northern BC). Report: (a) the current vacancy rate; (b) the average rent for a two-bedroom unit; (c) the annual change in each. Compare these to the corresponding figures for Vancouver. What do the differences tell you about relative supply and demand conditions?

\item \textbf{Supply-Demand Diagram.} Draw a supply-demand diagram illustrating the Canadian housing market. Then show the effect of each of the following, clearly labelling each shift as affecting supply or demand: (a) a 2-percentage-point increase in mortgage interest rates; (b) Canada increasing immigration to 500,000 per year; (c) widespread upzoning that allows four-unit buildings on all residential lots. For each, state the direction of the price and quantity effects.

\item \textbf{Affordability Calculation.} A household earns \$110,000 per year combined. The current 5-year fixed mortgage rate is 5.5\%. Under the stress test, they must qualify at 7.5\%.
  \begin{enumerate}[(a)]
    \item The standard affordability rule is that a mortgage can be at most 4.5$\times$ gross income under the stress test. What is the maximum purchase price if they put 20\% down?
    \item What down payment is required on a home priced at the national benchmark of \$730,000? If they put 5\% down, will CMHC insurance be required?
    \item Based on your answer, can this household realistically buy the national benchmark home? What about in Vancouver at \$1.2 million?
  \end{enumerate}

\item \textbf{Policy Essay (500 words).} Evaluate whether zoning reform alone can solve Canada's housing affordability crisis. Draw on the theoretical arguments, the international evidence (Auckland, Tokyo, Houston), and the specific challenges facing Canadian cities. Your essay should identify what zoning reform can and cannot accomplish, and propose what additional policies would be needed.

\end{enumerate}
